\documentclass[10pt, aspectratio=169]{beamer}

% Packages for XeTeX and table formatting
\usepackage{fontspec}
\usepackage{tabularx}
\usepackage{booktabs}
\usepackage{graphicx}
\usepackage{xcolor}
\usepackage{dirtree}
\usepackage[spanish]{babel}
% Theme selection
\usetheme{Madrid}
\usecolortheme{default}

% Presentation Information
\title[Participación Ciudadana]{
    Participación Ciudadana en la Ordenación Territorial y Urbanística: \\
    Comparativa entre la Legislación Estatal y Autonómica
}
\subject{Subject: Planeamiento Territorial y Urbanístico}
\author{Tomas González Llarena}
\date{\today}

\begin{document}

% --- Title Slide ---
\begin{frame}
    \titlepage
\end{frame}

% --- Frame 1 ---
\begin{frame}{1. Elaboración del planeamiento}
    \small
    \begin{tabularx}{\textwidth}{@{}X X@{}}
        \toprule
        \textbf{Legislación Estatal (Texto Refundido Ley de Suelo)} & \textbf{Legislación Autonómica (Ley 4/2017)} \\
        \midrule
        Se consagra el derecho de los ciudadanos a participar de forma efectiva en los procedimientos de elaboración y aprobación de cualquier instrumento de ordenación, así como en su evaluación ambiental. & Se reconoce el derecho a participar en los procedimientos de aprobación de los instrumentos de ordenación territorial, urbanística y de los recursos naturales. \\
        \midrule
        El derecho se materializa mediante la formulación de alegaciones, observaciones, propuestas, reclamaciones y quejas. & La participación se instrumenta a través de la formulación de alegaciones, observaciones y propuestas durante los periodos preceptivos de información pública. \\
        \midrule
        Se garantiza el derecho a obtener una respuesta motivada por parte de la Administración. & Se establece un proceso previo de participación ciudadana de carácter consultivo a través del portal web para instrumentos de iniciativa pública que requieran documento de avance. \\
        \bottomrule
    \end{tabularx}
 \end{frame}   

 

%  \begin{frame}{1.b Ejemplos de participación ciudadana en la elaboración del planeamiento}
%     % Input for external examples snippet
%     % \documentclass{article}

% \usepackage{xcolor} % Required for the colors

% \begin{document}

% The \dirtree command uses periods followed by a number to indicate depth.
% Always end each line with a period!
\dirtree{%
.1 \textbf{Participación ciudadana: Áreas Temáticas y Competencias}. %[cite: 6]
.2 \textbf{Movilidad y Accesibilidad}. %[cite: 6]
.3 \textcolor{magenta}{Competencia del Cabildo de Tenerife:}. %[cite: 6]
.4 Transporte público nocturno. %[cite: 6]
.4 Transporte público más barato para reducir el tráfico. %[cite: 6]
.4 No a la Vía Exterior. %[cite: 6]
.4 Guagua circular en el Suroeste. %[cite: 6]
.4 Aceras en TF-180 Vistabella. %[cite: 6]
.4 Miradores en Anaga. %[cite: 6]
.4 Mejora de la frecuencia del transporte público. %[cite: 6]
.4 Apartaderos carreteras Anaga. %[cite: 6]
.3 \textcolor{cyan}{Competencia del Ayuntamiento:}. %[cite: 6]
.4 Solución al estrechamiento carretera de acceso al Parque Rural por San Andrés. %[cite: 6]
.4 Conexión viaria entre barrios del Suroeste. %[cite: 6]
.4 Peatonalización de los centros urbanos. %[cite: 6]
.4 Ciudad para el peatón, no para el coche. %[cite: 6]
.4 Restringir aparcamientos en el centro SC (franja horaria). %[cite: 6]
.4 Determinar recorridos preferentes que sean accesibles a sillas de ruedas. %[cite: 6]
.2 \textbf{Espacio Libre Público}. %[cite: 6]
.3 \textcolor{cyan}{Competencia del Ayuntamiento:}. %[cite: 6]
.4 Uso de barrancos como zonas verdes. %[cite: 6]
.4 No a la privatización del espacio público (terrazas). %[cite: 6]
.4 Ejecución de los Espacios Libres previstos. %[cite: 6]
.4 Rutas deportivas iluminadas por la noche. %[cite: 6]
.4 Rutas deportivas segregadas y alejadas del tráfico. %[cite: 6]
.4 Renaturalización espacios urbanos. %[cite: 6]
.4 Apertura al mar en todo el municipio. %[cite: 6]
.2 \textbf{Dotaciones, Equipamientos e Infraestructuras}. %[cite: 6]
.3 \textcolor{cyan}{Competencia del Ayuntamiento:}. %[cite: 6]
.4 Pequeños intercambiadores modales en el perímetro de los barrios para dejar coche. %[cite: 6]
.4 Centros Dotacionales Integrales polivalentes. %[cite: 6]
.4 Centro de Salud en el Suroeste. %[cite: 6]
.4 Centros de participación cívica. %[cite: 6]
.4 {Ejecución dotaciones previstas en Acorán. Evitar ciudad dormitorio}. %[cite: 6]
.4 Mercadillo del Agricultor Suroeste. %[cite: 6]
.4 Camping entorno Las Teresitas. %[cite: 6]
.2 \textbf{Patrimonio Cultural}. %[cite: 6]
.3 \textcolor{cyan}{Competencia del Ayuntamiento:}. %[cite: 6]
.4 Protección del patrimonio inmaterial. %[cite: 6]
.4 Preservación de los estanques de agua para las aves migratorias y la agricultura. %[cite: 6]
.4 Revisar el Catálogo de Protección para ajustarlo a la realidad. %[cite: 6]
.4 Protección de eras, bancales, paseros. %[cite: 6]
.4 Rehabilitación del Barranco del Muerto. %[cite: 6]
.3 \textcolor{green}{Plan Estratégico Municipal:}. %[cite: 6]
.4 Fomentar comercio de proximidad. %[cite: 6]
.4 Ayudas a la Rehabilitación Edificios. %[cite: 6]
.4 Apostar por un parque de vivienda de alquiler. %[cite: 6]
.4 Divulgación del patrimonio material e inmaterial. %[cite: 6]
.2 \textbf{Desarrollo Sostenible}. %[cite: 6]
.3 \textcolor{cyan}{Competencia del Ayuntamiento:}. %[cite: 6]
.4 Sistemas bio urbanos, calles orientadas a corrientes de aire para evitar aires acondicionados. %[cite: 6]
.4 Rehabilitar el paisaje de la Montaña Taco. %[cite: 6]
.4 Creación de un parque agroecológico en el Suroeste. %[cite: 6]
.4 Huertos urbanos. %[cite: 6]
.4 Separar industrias molestas del área residencial. %[cite: 6]
.3 \textcolor{green}{Plan Estratégico Municipal:}. %[cite: 6]
.4 Lucha contra Cambio Climático. %[cite: 6]
.4 Fomento de las Energías Limpias. %[cite: 6]
.4 Hacer atractiva la ciudad para residir y no perder población. %[cite: 6]
}

% \end{document}
% \end{frame}

% --- Frame 2 ---
\begin{frame}{2. Información y transparencia}
    \begin{tabularx}{\textwidth}{@{}X X@{}}
        \toprule
        \textbf{Legislación Estatal (Texto Refundido Ley de Suelo)} & \textbf{Legislación Autonómica (Ley 4/2017)} \\
        \midrule
        Se reconoce el derecho a acceder a la información urbanística y ambiental, así como a ser informados por escrito del régimen aplicable a una finca. & Se garantiza el derecho a acceder a la información urbanística y ambiental, así como a ser informados por escrito del régimen aplicable a una finca. \\
        \midrule
        Se establece la obligación de las Administraciones de publicar en sus portales web los instrumentos de ordenación territorial y urbanística, así como los documentos de evaluación ambiental. & Se establece la obligación de las Administraciones de publicar en sus portales web los instrumentos de ordenación territorial y urbanística, así como los documentos de evaluación ambiental. \\
        \bottomrule
    \end{tabularx}
\end{frame}

% --- Frame 3 ---
\begin{frame}{3. Control de su aplicación}
    \begin{tabularx}{\textwidth}{@{}X X@{}}
        \toprule
        \textbf{Legislación Estatal (Texto Refundido Ley de Suelo)} & \textbf{Legislación Autonómica (Ley 4/2017)} \\
        \midrule
        Se reconoce el derecho a ejercer la acción pública con el fin de hacer respetar las determinaciones que indique la ordenación territorial y urbanística. & Se ampara el derecho a exigir el cumplimiento de la legalidad urbanística mediante la presentación de quejas y reclamaciones, y a través del ejercicio de la acción pública frente a órganos administrativos y judiciales. \\
        \midrule
        La acción pública también aplica para hacer respetar las decisiones resultantes de los procedimientos de evaluación ambiental de los proyectos. & Se impone a los particulares el deber de colaborar con el desarrollo de las funciones de control que la ley atribuye a las administraciones. \\
        \bottomrule
    \end{tabularx}
\end{frame}

% --- Frame 4 ---
\begin{frame}{4. Figuras de participación}
    \footnotesize % Slightly smaller font to fit the 5 rows
    \begin{tabularx}{\textwidth}{@{}X X@{}}
        \toprule
        \textbf{Legislación Estatal (Texto Refundido Ley de Suelo)} & \textbf{Legislación Autonómica (Ley 4/2017)} \\
        \midrule
        El derecho a la información urbanística y ambiental, permitiendo acceder a la información de las Administraciones y ser informados por escrito del régimen aplicable a una finca. & Las alegaciones, observaciones y propuestas en periodos de información pública. \\
        \midrule
        La presentación de propuestas, observaciones, quejas y alegaciones en los trámites de información pública. & Los procesos consultivos previos vía web. \\
        \midrule
        La acción pública en materia urbanística y medioambiental. & Los programas de participación ciudadana que pueden incluir sesiones abiertas al público, material divulgativo y consultas populares. \\
        \midrule
         & El ejercicio de la acción pública. \\
        \midrule
         & La figura del \"colaborador con el medioambiente y la sostenibilidad territorial\", que se adquiere mediante la firma de un convenio de colaboración con la Administración. \\
        \bottomrule
    \end{tabularx}
\end{frame}

\begin{frame}{Ejemplos de participación ciudadana en la elaboración del planeamiento}
    % \vspace{0.3cm}
    % Scale images heavily to prevent overflowing the slide
    \begin{center}
        \resizebox{!}{0.8\textheight}{
            \includegraphics{img/PPC_Res_Distritos.jpg}
        }
    \end{center}

 \end{frame}

\begin{frame}{Ejemplos de participación ciudadana en la elaboración del planeamiento (continuación)}
    % \vspace{0.3cm}
    % Scale images heavily to prevent overflowing the slide
    \begin{center}
       
        %\resizebox{!}{0.8\textheight}{
            \includegraphics[width=\textwidth, height=0.75\textheight, keepaspectratio]{img/ppc_1_movilidad.png}
        %    }
    \end{center}
\end{frame}

\begin{frame}{Ejemplos de participación ciudadana en la elaboración del planeamiento (continuación)}
    % \vspace{0.3cm}
    % Scale images heavily to prevent overflowing the slide
    \begin{center}
       
        %\resizebox{!}{0.8\textheight}{
            \includegraphics[width=\textwidth, height=0.75\textheight, keepaspectratio]{img/ppc_2_sostenible.png}
        %    }
    \end{center}
\end{frame}

\begin{frame}{Ejemplos de participación ciudadana en la elaboración del planeamiento (continuación)}
    % \vspace{0.3cm}
    % Scale images heavily to prevent overflowing the slide
    \begin{center}
       
        %\resizebox{!}{0.8\textheight}{
            \includegraphics[width=\textwidth, height=0.75\textheight, keepaspectratio]{img/ppc_3_ayuntamiento.png}
        %    }
    \end{center}
\end{frame}

\begin{frame}{Ejemplos de participación ciudadana en la elaboración del planeamiento (continuación)}
    % \vspace{0.3cm}
    % Scale images heavily to prevent overflowing the slide
    \begin{center}
       
        %\resizebox{!}{0.8\textheight}{
            \includegraphics[width=\textwidth, height=0.75\textheight, keepaspectratio]{img/ppc_4_cabildo.png}
        %    }
    \end{center}
\end{frame}

\begin{frame}{Ejemplos de participación ciudadana en la elaboración del planeamiento (continuación)}
    % \vspace{0.3cm}
    % Scale images heavily to prevent overflowing the slide
    \begin{center}
       
        %\resizebox{!}{0.8\textheight}{
            \includegraphics[width=\textwidth, height=0.75\textheight, keepaspectratio]{img/ppc_5_espacio_libre.png}
        %    }
    \end{center}
\end{frame}

\begin{frame}{Ejemplos de participación ciudadana en la elaboración del planeamiento (continuación)}
    % \vspace{0.3cm}
    % Scale images heavily to prevent overflowing the slide
    \begin{center}
       
        %\resizebox{!}{0.8\textheight}{
            \includegraphics[width=\textwidth, height=0.75\textheight, keepaspectratio]{img/ppc_6_patrimonio.png}
        %    }
    \end{center}
\end{frame}

\begin{frame}{Ejemplos de participación ciudadana en la elaboración del planeamiento (continuación)}
    % \vspace{0.3cm}
    % Scale images heavily to prevent overflowing the slide
    \begin{center}
       
        %\resizebox{!}{0.8\textheight}{
            \includegraphics[width=\textwidth, height=0.75\textheight, keepaspectratio]{img/ppc_7_PEM.png}
        %    }
    \end{center}
\end{frame}


\begin{frame}
    \frametitle{Conclusión}
    La legilacion autonómica (Ley 4/2017) hace un esfuerzo {\it infinitesimalmente pequeño} para establecer un marco más amplio y detallado para la participación ciudadana en la ordenación territorial y urbanística, 
    incluyendo:
    \begin{itemize}
        \item procesos consultivos previos,
        \item programas de participación ciudadana
        \item la figura del colaborador con el medioambiente y la sostenibilidad territorial.
    \end{itemize}
    Esto refleja un enfoque más inclusivo y participativo en comparación con la legislación estatal (Texto Refundido Ley de Suelo), que se centra principalmente en:
    \begin{itemize}
        \item el derecho a la información
        \item la presentación de alegaciones y
        \item la acción pública.
    \end{itemize}
\end{frame}

\end{document}